\documentclass[10pt,a4paper]{article}
\usepackage[utf8]{inputenc}
\usepackage[T1]{fontenc}
\usepackage[margin=1.8cm,top=2.6cm,bottom=2.2cm]{geometry}
\usepackage{parskip}
\usepackage{microtype}
\usepackage{amsmath,amssymb}
\usepackage{array,tabularx,multirow}
\usepackage[table]{xcolor}
\usepackage{graphicx,float}
\usepackage{titlesec}
\usepackage{enumitem}
\usepackage{fancyhdr}

\titleformat{\section}{\large\bfseries}{}{0em}{}[\titlerule]
\titleformat{\subsection}{\normalsize\bfseries}{}{0em}{}
\titleformat{\subsubsection}{\normalsize\itshape}{}{0em}{}
\titlespacing{\section}{0pt}{1.4ex plus .2ex}{0.5ex}
\titlespacing{\subsection}{0pt}{0.9ex}{0.3ex}

% Page headers
\pagestyle{fancy}
\fancyhf{}
\setlength{\headheight}{15pt}
\fancyhead[L]{\small\textit{Layer~0: Wound Biology}}
\fancyhead[R]{\small Aaran Nihalani $\bullet$ Eton iGEM Dry Lab}
\fancyfoot[C]{\small\thepage}
\renewcommand{\headrulewidth}{0.35pt}

% Colour palette -- muted
\definecolor{boxblue}{rgb}{0.93,0.96,0.99}
\definecolor{boxgreen}{rgb}{0.91,0.97,0.92}
\definecolor{boxamber}{rgb}{0.99,0.97,0.88}
\definecolor{boxgray}{rgb}{0.95,0.95,0.95}
\definecolor{boxred}{rgb}{0.99,0.92,0.92}
\definecolor{ambertext}{rgb}{0.46,0.24,0.00}
\definecolor{alertred}{rgb}{0.68,0.06,0.06}

\setlength{\fboxsep}{7pt}
\newsavebox{\tmpbox}

\newenvironment{infobox}{%
  \begin{lrbox}{\tmpbox}\begin{minipage}{0.965\textwidth}\vspace{0.2em}}%
 {\vspace{0.2em}\end{minipage}\end{lrbox}%
  \par\vspace{0.4em}\noindent\colorbox{boxblue}{\usebox{\tmpbox}}\par\vspace{0.4em}}

\newenvironment{findingbox}{%
  \begin{lrbox}{\tmpbox}\begin{minipage}{0.965\textwidth}\vspace{0.2em}}%
 {\vspace{0.2em}\end{minipage}\end{lrbox}%
  \par\vspace{0.4em}\noindent\colorbox{boxgreen}{\usebox{\tmpbox}}\par\vspace{0.4em}}

\newenvironment{needsbox}{%
  \begin{lrbox}{\tmpbox}\begin{minipage}{0.965\textwidth}\vspace{0.2em}%
  \textbf{\textcolor{ambertext}{Outstanding measurements:}}\par\vspace{0.15em}}%
 {\vspace{0.2em}\end{minipage}\end{lrbox}%
  \par\vspace{0.4em}\noindent\colorbox{boxamber}{\usebox{\tmpbox}}\par\vspace{0.4em}}

\newenvironment{alertbox}{%
  \begin{lrbox}{\tmpbox}\begin{minipage}{0.965\textwidth}\vspace{0.2em}%
  \textbf{\textcolor{alertred}{Critical finding:}}\par\vspace{0.15em}}%
 {\vspace{0.2em}\end{minipage}\end{lrbox}%
  \par\vspace{0.4em}\noindent\colorbox{boxred}{\usebox{\tmpbox}}\par\vspace{0.4em}}

\newenvironment{eqbox}{%
  \begin{lrbox}{\tmpbox}\begin{minipage}{0.965\textwidth}\vspace{0.2em}}%
 {\vspace{0.2em}\end{minipage}\end{lrbox}%
  \par\vspace{0.3em}\noindent\colorbox{boxgray}{\usebox{\tmpbox}}\par\vspace{0.3em}}

% Macros
\newcommand{\litref}[1]{\textsuperscript{\normalfont[#1]}}
\newcommand{\kon}{k_{\text{on}}}
\newcommand{\koff}{k_{\text{off}}}

% Table status symbols
\newcommand{\symconf}{$\checkmark$}     % confirmed from paper text
\newcommand{\symapprox}{${\sim}$}       % approximate from figure
\newcommand{\symneed}{$\circ$}          % needs citation/measurement

\begin{document}

%----------------------------------------------------------------------
\begin{center}
{\LARGE\bfseries Layer 0: Wound Biology}\\[0.4em]
{\large Protease Cascade ODE: MMP-9 and HNE in Chronic Wound Fluid}\\[0.25em]
{\small\color{gray} Eton iGEM 2026, Dry Lab \quad|\quad
Layer 0 of 17-Layer Biosensor Modelling Pipeline \quad|\quad
\texttt{wound\_biology.py}}
\end{center}
\vspace{0.4em}\hrule\vspace{0.7em}

%----------------------------------------------------------------------
\section{Purpose}

Layer~0 predicts the \textbf{free protease concentration in wound fluid at the time of
swab sampling} and how it evolves over the 2-hour post-sampling window.
Its outputs, $[\text{MMP-9}]_{\text{free}}(t)$ and $[\text{HNE}]_{\text{free}}(t)$,
are in-situ wound fluid concentrations. Layer~1 (sample collection) applies
the swab dilution factor $D_f$ to convert from in-situ concentrations to what enters the device.

The distinction between healing and chronic wound states is encoded through two distinct
initial-condition parameter sets, both grounded in published ELISA data.
The model is a system of ODEs for nine coupled species: pro-MMP-9, free MMP-9, free TIMP-1,
MMP-9:TIMP-1 complex, free HNE, free A1AT, HNE:A1AT complex, free SLPI, and HNE:SLPI complex.

\begin{infobox}
\textbf{Parameter citation policy.} Every parameter used by this model must be grounded
in a paper in \texttt{electrochemical\_device/papers/layer\_00/} (status: \symconf),
approximate from a figure in such a paper (status: \symapprox), or explicitly flagged
as awaiting measurement (status: \symneed). No values from secondary literature alone.
\end{infobox}

%----------------------------------------------------------------------
\section{I. \; Wound Fluid Composition}

\subsection{Biological rationale}

Chronic wounds are characterised by a prolonged inflammatory phase driven by persistent
neutrophil and macrophage infiltration\litref{1}. This sustains elevated protease secretion
long after acute wounds have resolved. The two proteases targeted by this device are:

\begin{itemize}[noitemsep]
  \item \textbf{MMP-9 (gelatinase B)}: secreted by neutrophils and macrophages as an
    inactive zymogen (pro-MMP-9). Activation requires proteolytic removal of the prodomain.
    In chronic wounds, both forms are elevated; the ratio of active to pro form predicts
    healing outcome.
  \item \textbf{HNE (human neutrophil elastase)}: a serine protease released already
    active from azurophil granules. Regulated primarily by $\alpha_1$-antitrypsin (A1AT)
    and secondarily by SLPI.
\end{itemize}

\subsection{Literature evidence}

\textbf{Trengove et al.\ (1999)\litref{1}} measured protease activity in wound fluid
from chronic venous leg ulcers (VLU) versus acute wounds (all values from paper text, p.446--448):
\begin{itemize}[noitemsep]
  \item Chronic VLU MMP activity: median $22.8\;\mu\text{g/mL}$, versus $0.76\;\mu\text{g/mL}$
    in acute wounds ($\approx30\times$ elevation, $p<0.001$; 92-kDa MMP-9 predominant on zymography).
  \item TIMP-1 (ELISA): chronic VLU median $0.8\;\mu\text{g/mL} = 800$~ng/mL ($n=5$; text p.448).
  \item HNE (activity assay): non-healing VLU median $1.4\;\mu\text{g/mL} \approx 47$~nM
    (range $<1$--$11.9\;\mu\text{g/mL}$; elevated in only 6 of 13 chronic samples; text p.448).
  \item MMP activity decreased significantly as wounds healed ($p < 0.01$; text p.449).
\end{itemize}

\textbf{Ladwig et al.\ (2002)\litref{2}} used quantitative gelatin zymography and ELISA in
wound fluid from 56 patients with chronic pressure ulcers. Key values (all from paper text p.31):
\begin{itemize}[noitemsep]
  \item Pro-MMP-9: mean $8{,}100 \pm 3{,}000$~ng/mL.
    Active MMP-9: mean $4{,}800 \pm 1{,}300$~ng/mL.
  \item Total MMP-9 $= 8{,}100 + 4{,}800 = 12{,}900$~ng/mL; pro-form fraction
    $= 8{,}100/12{,}900 = 63\%$.
  \item Active MMP-9 significantly higher in poor healers at Day~0 (text p.33;
    approximate values from Figure~8).
  \item TIMP-1 significantly higher in good healers at Day~0 (text p.33;
    approximate values from Figure~10).
  \item The MMP-9:TIMP-1 ratio was a predictor of healing trajectory (text p.34, Figure~11).
\end{itemize}

\textbf{Rao et al.\ (1995)\litref{7}} examined A1AT integrity in wound fluid from chronic
and healing wounds (J Invest Dermatol 105:572--578, text p.572--575). In healing and acute
wound fluid, A1AT was intact and functional in all samples. In chronic wound fluid, A1AT was
degraded and non-functional in 9 of 10 samples; the authors attribute this to sustained
serine and metalloprotease activity in the chronic wound environment. This finding is critical
for the HNE channel: functional A1AT cannot be assumed present in chronic wounds.

\textbf{Nwomeh et al.\ (1999)\litref{8}} measured matrix metalloproteinase activities and
inhibitor levels in venous leg ulcers (J Surg Res 81:189--195). TIMP-1 levels in chronic
ulcer fluid were ``very low or not detected'' (text p.193), consistent with the low TIMP-1
values in Trengove 1999\litref{1} and confirming TIMP-1 suppression as a hallmark of the
chronic wound environment.

\begin{needsbox}
\textbf{Healing wound concentrations (not confirmed):} Ladwig 2002 is a pressure ulcer
study; the healing scenario MMP-9 ($\approx5{,}400$~ng/mL) is read approximately from
Figure~8 (not stated in text). The healing TIMP-1 ($= 1{,}104$~ng/mL) has no
paper-text source -- it is a working estimate. The healing wound prediction (Section~V)
is therefore illustrative. It depends on the sign of the MMP-9/TIMP-1 excess, not the
exact values; a confirmed joint VLU ELISA (MMP-9 + TIMP-1 in the same cohort, chronic
and healing time points) is needed before the discrimination ratio is finalised.
\end{needsbox}

\subsection{Initial conditions}

Status: \symconf~confirmed from paper text, \symapprox~approximate from figure,
\symneed~needs measurement.

\begin{center}
\setlength{\tabcolsep}{6pt}
\renewcommand{\arraystretch}{1.35}
\begin{tabular}{>{\small}p{3.8cm}
                >{\small\raggedright\arraybackslash}p{3.3cm}
                >{\small\raggedright\arraybackslash}p{3.3cm}
                >{\small\centering\arraybackslash}p{0.8cm}
                >{\small\raggedright\arraybackslash}p{2.8cm}}
\hline\hline
\textbf{Parameter} & \textbf{Chronic wound} & \textbf{Healing wound} & \textbf{St.} & \textbf{Source} \\\hline
$[\text{MMP-9}]_{\rm total}$
  & 12{,}900~ng/mL (140~nM) & 5{,}400~ng/mL (59~nM)
  & \symconf/\symapprox & [2] text/Fig.~8 \\
$[\text{pro-MMP-9}]$ (63\%)
  & 8{,}100~ng/mL (88~nM) & 3{,}400~ng/mL (37~nM)
  & \symconf/\symapprox & [2] text \\
$[\text{active MMP-9}]$ (37\%)
  & 4{,}800~ng/mL (52~nM) & 2{,}000~ng/mL (22~nM)
  & \symconf/\symapprox & [2] text/Fig.~8 \\
\hline
$[\text{TIMP-1}]$
  & 800~ng/mL (28~nM) & 1{,}104~ng/mL (39~nM)
  & \symconf/\symneed & [1] text/estimate \\
\hline
$[\text{HNE}]$ (both states)
  & \multicolumn{2}{c}{47~nM (1.4~$\mu$g/mL)}
  & \symconf & [1] text \\
$[\text{A1AT}]$ (functional)
  & $\approx 0$ (degraded) & 1.92~$\mu$M (estimate)
  & \symconf/\symneed & [7] text/est. \\\hline\hline
\end{tabular}
\end{center}

\textbf{Physical basis for wound-state discrimination.}
In chronic wounds, active MMP-9 (52~nM) exceeds TIMP-1 (28~nM) by 24~nM. Since
$K_i(\text{TIMP-1/MMP-9}) = 8.5$~nM (Olson 1997\litref{4}, confirmed), which is small
compared to both concentrations, TIMP-1 is nearly fully bound and free MMP-9 persists at
$\approx 32$~nM (ODE result, t = 30~min). In healing wounds (estimated), TIMP-1 (39~nM)
exceeds active MMP-9 (22~nM) by 17~nM, suppressing free MMP-9 to $\approx 6$~nM.

A1AT integrity provides a second discriminating axis for HNE (Rao 1995\litref{7}): in healing
wounds A1AT is intact, quenching HNE to 0.00~pM within milliseconds; in chronic wounds A1AT
is degraded in 9/10 samples, leaving HNE partially regulated only by SLPI ($\approx$30~nM
at $t = 30$~min). However, elevated HNE was detected in only 6 of 13 chronic VLU samples
(Trengove 1999\litref{1}), limiting its reliability as a primary discriminator.

%----------------------------------------------------------------------
\section{II. \; Activation Kinetics (MMP-9 Pathway)}

Pro-MMP-9 activation by MMP-3 and plasmin is modelled as pseudo-first-order in pro-MMP-9:

\begin{eqbox}
\begin{equation}
  \frac{d[\text{pro-MMP-9}]}{dt}
  = -\left(\frac{k_{\rm cat}}{K_m}\right)_{\rm MMP3}\![\text{MMP-3}]\,[\text{pro-MMP-9}]
    -\left(\frac{k_{\rm cat}}{K_m}\right)_{\rm plasmin}\![\text{plasmin}]\,[\text{pro-MMP-9}]
\end{equation}
\end{eqbox}

\begin{needsbox}
$(k_{\rm cat}/K_m)_{\rm MMP3}$, $(k_{\rm cat}/K_m)_{\rm plasmin}$, $[\text{MMP-3}]$, and
$[\text{plasmin}]$ in wound fluid are all unconfirmed placeholders. No paper in the papers
folder provides wound-fluid MMP-3 or plasmin concentrations or the relevant kinetic constants.
With current placeholder values, $<1$\% of pro-MMP-9 is activated in 120~min, so these
values have no effect on the chronic wound free MMP-9 prediction. They would matter only if
pro-MMP-9 activation within the 2-hour window were a significant fraction of the total.
\end{needsbox}

%----------------------------------------------------------------------
\section{III. \; Inhibition Kinetics}

\subsection{TIMP-1 / MMP-9 (reversible tight inhibition)}

TIMP-1 binds to the MMP-9 active site in a non-covalent, reversible manner with biphasic
kinetics (two binding sites; high-affinity site dominates at wound fluid concentrations):

\begin{eqbox}
\begin{equation}
  \text{MMP-9} + \text{TIMP-1} \underset{\koff}{\overset{\kon}{\rightleftharpoons}}
  \text{MMP-9:TIMP-1}
  \qquad K_i = \frac{\koff}{\kon}
\end{equation}
\end{eqbox}

\textbf{Confirmed kinetic constants (Olson et al.\ 1997\litref{4}, Table~IV,
enzyme inhibition assay, MMP-9 82~kDa active form):}

\begin{eqbox}
\begin{align*}
  \kon &= 2.48 \pm 0.47 \times 10^5\;\text{M}^{-1}\text{s}^{-1} \\
  \koff &= 2.1 \pm 1.1 \times 10^{-3}\;\text{s}^{-1} \\
  K_i &= \koff/\kon = 8.5 \pm 2.9\;\text{nM}
    \quad\text{(stated directly in Table~IV; check: }2.1\times10^{-3}/2.48\times10^5 = 8.47\;\text{nM)}
\end{align*}
\end{eqbox}

Enzyme inhibition assay values are used in preference to SPR values (Olson 1997
Tables~I and II) because they reflect solution-phase kinetics at conditions closer to
wound fluid. As noted by Olson 1997 (p.29981), enzyme inhibition $\kon$ values are
consistently 2.5--7$\times$ higher than SPR values, making them more appropriate for
modelling diffusion-limited association in wound fluid.

\textbf{Sensitivity of discrimination ratio to $K_i$.}
The chronic wound prediction ($\approx32$~nM free MMP-9 at 30~min) is robust: changing
$K_i$ from 8.5~nM to 1~nM shifts the result from 32~nM to 27~nM (a 15\% change).
The healing prediction is more sensitive: $K_i = 8.5$~nM gives 6~nM free MMP-9,
whereas $K_i = 1$~nM gives 1.3~nM. This is because the healing wound MMP-9/TIMP-1 excess
is only 17~nM, comparable to $K_i$.

\subsection{A1AT / HNE (near-irreversible inhibition)}

A1AT forms a covalent acyl-enzyme intermediate with HNE that is effectively irreversible:

\begin{eqbox}
\begin{equation}
  \text{HNE} + \text{A1AT} \xrightarrow{\kon} \text{HNE:A1AT}
  \quad(\text{irreversible; }\koff \approx 0)
\end{equation}
\end{eqbox}

\textbf{Confirmed kinetic constant (Beatty, Bieth \& Travis 1980\litref{5}, Table~1):}
\begin{eqbox}
\begin{equation}
  \kon(\text{HNE}/\text{A1AT}) = 6.5 \pm 4.0 \times 10^7\;\text{M}^{-1}\text{s}^{-1}
  \quad\text{(at 25\,\textdegree C, pH 8.0)}
\end{equation}
\end{eqbox}

\textbf{Half-life of free HNE in healing wounds} (with functional $[\text{A1AT}]_0 = 1.92$~$\mu$M, estimated):
\begin{eqbox}
\begin{equation}
  t_{1/2}^{\rm HNE}
  = \frac{\ln 2}{\kon \cdot [\text{A1AT}]_0}
  = \frac{0.693}{6.5\times10^7 \times 1.92\times10^{-6}}
  = 0.0055\;\text{s} \quad (5.5\;\text{ms})
\end{equation}
\end{eqbox}

Because $[\text{A1AT}]_0$ is an estimate, the conclusion is verified at 10$\times$ lower
$[\text{A1AT}]_0$: $t_{1/2} = 55$~ms, still three orders of magnitude below the
5-minute measurement window. This analysis applies to healing wounds where A1AT is intact
(Rao et al.\ 1995\litref{7}). In chronic wounds, A1AT is degraded in 9 of 10 samples;
see the finding below.

\begin{alertbox}
\textbf{A1AT is wound-state dependent (Rao et al.\ 1995\litref{7}).}
In \textbf{healing} wounds A1AT is intact; HNE is quenched to 0.00~pM within 5.5~ms
(confirmed $\kon = 6.5\times10^7$~M$^{-1}$s$^{-1}$, Beatty 1980\litref{5}; insensitive
to [A1AT] estimate: $t_{1/2} = 55$~ms even at 10$\times$ lower A1AT). In \textbf{chronic}
wounds, A1AT is degraded and non-functional in 9/10 samples (Rao 1995\litref{7}); HNE is
only partially inhibited by SLPI, persisting at $\approx$30~nM at $t = 30$~min.

\textbf{HNE is an unreliable primary sensor target.} Although HNE discriminates between
wound states (chronic $\approx$30~nM vs healing 0~pM), elevated HNE was detected in only
6 of 13 chronic VLU samples (Trengove 1999\litref{1}): the remaining 50\% of chronic
cases have HNE at or below detection. MMP-9 is the preferred primary target. The second
electrode should serve as an internal reference or control, not a primary HNE sensing channel.
\end{alertbox}

\subsection{SLPI / HNE (secondary inhibitor)}

SLPI is a serine protease inhibitor expressed in wound healing tissue\litref{6}.
Placeholder kinetic constants ($\kon = 10^5$~M$^{-1}$s$^{-1}$, $\koff = 10^{-4}$~s$^{-1}$,
concentration $\approx$200~ng/mL $\approx$ 17~nM) are used. Ashcroft et al.\ (2000)\litref{6}
confirmed the importance of SLPI to wound healing outcomes in knockout mice but did not
state a quantitative SLPI concentration in human wound fluid.

In healing wounds, HNE is completely neutralised by intact A1AT (Rao 1995\litref{7}),
making SLPI irrelevant. In chronic wounds, A1AT is degraded in 9/10 samples (Rao
1995\litref{7}), and SLPI provides the primary remaining HNE inhibition; the ODE predicts
$\approx$30~nM free HNE at $t = 30$~min after partial SLPI inhibition. The chronic wound
HNE prediction is sensitive to placeholder SLPI kinetics and concentration, which remain
unconfirmed from papers in this folder.

%----------------------------------------------------------------------
\section{IV. \; Full ODE System}

\noindent Variable key (all concentrations in mol/L):

\vspace{3pt}
\begin{center}
\small
\renewcommand{\arraystretch}{1.15}
\begin{tabular}{@{}ll@{\qquad}ll@{}}
$P$ & pro-MMP-9           & $C_1$ & MMP-9:TIMP-1 complex \\
$M$ & free (active) MMP-9 & $H$   & free HNE \\
$T$ & free TIMP-1         & $A$   & free A1AT \\
$C_2$ & HNE:A1AT complex  & $S$   & free SLPI \\
$C_3$ & HNE:SLPI complex  &       & \\
\end{tabular}
\end{center}

\begin{eqbox}
\small
\begin{align}
  \dot{P}   &= -k^*\,P \\[2pt]
  \dot{M}   &= +k^*\,P  - \kon^{T} MT + \koff^{T} C_1 \\[2pt]
  \dot{T}   &= -\kon^{T} MT + \koff^{T} C_1 \\[2pt]
  \dot{C}_1 &= +\kon^{T} MT - \koff^{T} C_1 \\[2pt]
  \dot{H}   &= -\kon^{A} HA - \kon^{S} HS + \koff^{S} C_3 \\[2pt]
  \dot{A}   &= -\kon^{A} HA \\[2pt]
  \dot{C}_2 &= +\kon^{A} HA \\[2pt]
  \dot{S}   &= -\kon^{S} HS + \koff^{S} C_3 \\[2pt]
  \dot{C}_3 &= +\kon^{S} HS - \koff^{S} C_3
\end{align}
\vspace{-4pt}
\noindent where $k^* = (k_{\rm cat}/K_m)_{\rm MMP3}[{\rm MMP\text{-}3}]
+ (k_{\rm cat}/K_m)_{\rm plasmin}[{\rm plasmin}]$;\;
$k_{\rm on,off}^{T}$ for TIMP-1/MMP-9 (Olson 1997\litref{4});\;
$\kon^{A}$ for A1AT/HNE (Beatty 1980\litref{5}).
\end{eqbox}

Solved with the Radau stiff solver (\texttt{scipy.integrate.solve\_ivp},
\texttt{rtol=1e-10, atol=1e-15}). Stiff integration is required because the A1AT/HNE
timescale ($\sim$5.5~ms) is $10^7$-fold faster than the pro-MMP-9 activation timescale
($\sim$30--120~min). The TIMP-1/MMP-9 equilibration is intermediate ($\tau \approx 2$~min).

%----------------------------------------------------------------------
\section{V. \; Model Output and Wound-State Discrimination}

\subsection{Discrimination between wound states at $t = 30$ min}

\begin{center}
\setlength{\tabcolsep}{6pt}
\renewcommand{\arraystretch}{1.35}
\begin{tabular}{>{\small}p{4.4cm}
                >{\small\raggedright\arraybackslash}p{3.0cm}
                >{\small\raggedright\arraybackslash}p{3.0cm}
                >{\small\centering\arraybackslash}p{1.1cm}
                >{\small\raggedright\arraybackslash}p{2.0cm}}
\hline\hline
\textbf{Parameter} & \textbf{Chronic wound} & \textbf{Healing wound} & \textbf{C/H} & \textbf{Source} \\\hline
$[\text{MMP-9}]_{\rm total}$ (in-situ)
  & 12{,}900~ng/mL (140~nM) & 5{,}400~ng/mL (59~nM)
  & $2.4\times$ & [2] text/Fig. \\
$[\text{MMP-9}]_{\rm active}$ (in-situ)
  & 4{,}800~ng/mL (52~nM) & 2{,}000~ng/mL (22~nM)
  & $2.4\times$ & [2] text/Fig. \\
$[\text{TIMP-1}]$ (in-situ)
  & 800~ng/mL (28~nM) & 1{,}104~ng/mL (39~nM)
  & $0.72\times$ & [1] text/est. \\
MMP-9 vs TIMP-1 excess
  & $+24$~nM (MMP-9) & $-17$~nM (TIMP-1)
  & -- & derived \\
\hline
\multicolumn{5}{l}{\footnotesize\textit{ODE output at $t = 30$~min (in-situ, before Layer~1 dilution $D_f$):}} \\[1pt]
$[\text{MMP-9}]_{\rm free}$
  & $\mathbf{32}$~nM & $\mathbf{6}$~nM
  & $\mathbf{5\times}$ & ODE \\
$[\text{TIMP-1}]_{\rm free}$
  & $\approx$6~nM & $\approx$22~nM
  & -- & ODE \\
$[\text{HNE}]_{\rm free}$
  & $\approx$30~nM$^{\dagger}$ & 0.00~pM
  & -- & ODE\litref{7} \\\hline\hline
\end{tabular}
\end{center}
{\footnotesize $\dagger$~A1AT degraded in chronic wounds (Rao 1995\litref{7});
HNE elevated in only 50\% of chronic VLU (Trengove 1999\litref{1}).
Chronic HNE $\approx$30~nM is sensitive to placeholder SLPI kinetics.}

\begin{findingbox}
\textbf{Confirmed chronic wound prediction.}
With Ladwig 2002 text-confirmed concentrations (active MMP-9 $= 52$~nM, TIMP-1 $= 28$~nM)\litref{2}
and Olson 1997 confirmed kinetics ($K_i = 8.5$~nM)\litref{4}, the ODE gives free MMP-9
$\approx 32$~nM at $t = 30$~min in chronic wounds. This is driven entirely by the
24~nM MMP-9 excess over TIMP-1; no production rate assumptions are needed.

\textbf{Illustrative healing comparison.}
The healing result ($\approx 6$~nM) depends on unconfirmed estimates: healing MMP-9
from Ladwig Figure~8 (approximate) and healing TIMP-1 with no paper source. The
discrimination ratio ($\sim$5$\times$) should be treated as a working estimate until
a joint VLU ELISA is performed. The mechanism is sound: when TIMP-1 excess exceeds
active MMP-9, free MMP-9 is suppressed; the exact ratio follows from the confirmed $K_i$.

\textbf{HNE secondary channel.}
Confirmed A1AT wound-state dependence (Rao 1995\litref{7}) gives strong HNE discrimination:
chronic $\approx$30~nM vs healing 0.00~pM. However, elevated HNE was observed in only
6/13 (50\%) of chronic VLU samples (Trengove 1999\litref{1}); the other 50\% show HNE
at or below detection. HNE is therefore not a reliable standalone discriminator and the
30~nM chronic estimate is sensitive to unconfirmed SLPI kinetics. MMP-9 ($\sim$5$\times$
discrimination, grounded in confirmed kinetics) is the preferred primary target.
\end{findingbox}

\subsection{Layer~1 dilution and sensor range}

For chronic wound free MMP-9 (32~nM) to land in the sensor range (17--200~ng/mL $= 0.185$--2.17~nM):

\begin{eqbox}
\begin{equation}
  D_f = \frac{[\text{MMP-9}]_{\rm free,\,in\text{-}situ}}{[\text{MMP-9}]_{\rm target}}
  = \frac{32\;\text{nM}}{0.185\text{--}2.17\;\text{nM}} \approx 15\text{--}173
\end{equation}
\end{eqbox}

A dilution factor $D_f \approx 10$--100 is within the range of practical swab designs.
Example with $D_f = 50$: chronic wound device entry $= 32/50 = 0.64$~nM $= 59$~ng/mL
(within sensor range); healing wound $= 6/50 = 0.12$~nM $= 11$~ng/mL (near the lower end).
The healing signal approaching the LOD under large $D_f$ is a risk to address in Layer~1.

%----------------------------------------------------------------------
\section{VI. \; Parameter Summary}

\textbf{Status}: \symconf~confirmed from paper text, \symapprox~from figure (not text), \symneed~needs measurement/citation.
\textbf{Sensitivity}: whether changing the value by 10$\times$ affects the
device-relevant predictions.

\begin{center}
{\small
\setlength{\tabcolsep}{4pt}
\renewcommand{\arraystretch}{1.28}
\begin{tabular}{p{3.0cm}p{2.2cm}p{4.4cm}cp{1.5cm}p{1.2cm}}
\hline\hline
\textbf{Parameter} & \textbf{Value} & \textbf{Source} & \textbf{St.} & \textbf{Affects} & \textbf{Sensitv.} \\\hline
$[\text{MMP-9}]_{\rm pro}$ (chr.)
  & 8{,}100~ng/mL
  & Ladwig 2002\litref{2} text p.31
  & \symconf & I.C. & Low \\
$[\text{MMP-9}]_{\rm act}$ (chr.)
  & 4{,}800~ng/mL
  & Ladwig 2002\litref{2} text p.31
  & \symconf & I.C. & High \\
Pro-form fraction
  & 63\%
  & Ladwig 2002\litref{2} (8100/12900)
  & \symconf & I.C. & Low \\
$[\text{TIMP-1}]$ (chronic)
  & 800~ng/mL
  & Trengove 1999\litref{1} text p.448
  & \symconf & MMP-9 eq. & High \\
$[\text{HNE}]$
  & 47~nM
  & Trengove 1999\litref{1} text p.448
  & \symconf & HNE (n/v) & n/a \\
$[\text{MMP-9}]_{\rm act}$ (heal.)
  & 2{,}000~ng/mL
  & Ladwig 2002\litref{2} Fig.~8 (approx.)
  & \symapprox & I.C. (heal.) & High \\
$[\text{TIMP-1}]$ (healing)
  & 1{,}104~ng/mL
  & No paper in folder -- estimate
  & \symneed & MMP-9 eq. & High \\
$\kon$ (TIMP-1/MMP-9)
  & $2.48\!\times\!10^5$~M$^{-1}$s$^{-1}$
  & Olson 1997\litref{4} Table IV
  & \symconf & MMP-9 eq. & Low \\
$\koff$ (TIMP-1/MMP-9)
  & $2.1\!\times\!10^{-3}$~s$^{-1}$
  & Olson 1997\litref{4} Table IV
  & \symconf & $K_i$ & Moderate \\
$K_i$ (TIMP-1/MMP-9)
  & 8.5~nM
  & Olson 1997\litref{4} Table IV
  & \symconf & MMP-9 eq. & Moderate \\
$\kon$ (A1AT/HNE)
  & $6.5\!\times\!10^7$~M$^{-1}$s$^{-1}$
  & Beatty 1980\litref{5} Table~1
  & \symconf & HNE (heal.) & Low \\
$[\text{A1AT}]$ (healing)
  & 1.92~$\mu$M (est.)
  & Serum $\div$10; intact (Rao 1995\litref{7})
  & \symconf/\symneed & HNE 0~pM & Low* \\
$[\text{A1AT}]$ (chronic)
  & $\approx 0$ (degraded)
  & Rao 1995\litref{7} text: 9/10 samples
  & \symconf & HNE $\approx$30~nM & High \\
$[\text{MMP-3}]$
  & 10~ng/mL
  & Placeholder
  & \symneed & Act. rate & None \\
$[\text{plasmin}]$
  & 1~nM
  & Placeholder
  & \symneed & Act. rate & None \\
$k_{\rm cat}/K_m$ (MMP-3)
  & $5\!\times\!10^3$~M$^{-1}$s$^{-1}$
  & Placeholder
  & \symneed & Act. rate & None \\
$[\text{SLPI}]$
  & 200~ng/mL
  & Ashcroft 2000\litref{6}: no conc.\ stated
  & \symneed & HNE (chr.) & Moderate \\\hline\hline
\end{tabular}
}
\end{center}

\textit{\small *In healing wounds, A1AT intact (Rao 1995\litref{7}): HNE $t_{1/2} = 55$~ms
even at 10$\times$ lower [A1AT]; healing HNE $= 0.00$~pM is insensitive to this estimate.
In chronic wounds, A1AT $\approx 0$ (confirmed degraded); chronic HNE prediction depends on
placeholder SLPI kinetics (Moderate sensitivity).}

%----------------------------------------------------------------------
\section{References}

All papers in \texttt{electrochemical\_device/papers/layer\_00/}
unless otherwise noted.

\begin{enumerate}[label={[\arabic*]},noitemsep]
\item Trengove NJ, Stacey MC, Macauley S et al.\ (1999) Analysis of the acute and chronic
  wound environments: the role of proteases and their inhibitors.
  \textit{Wound Repair Regen} 7(6):442--52.
  \textbf{trengove1999.pdf} -- values from text confirmed.

\item Ladwig GP, Robson MC, Liu R et al.\ (2002) Ratios of activated matrix
  metalloproteinase-9 to tissue inhibitor of metalloproteinase-1 in wound fluids
  are inversely correlated with healing of pressure ulcers.
  \textit{Wound Repair Regen} 10(1):26--37.
  \textbf{ladwig2002.pdf} -- chronic concentrations from text (p.31) confirmed.
  Healing concentrations from Figure~8 only (approximate).

\item Nagase H, Woessner JF~Jr (1999) Matrix metalloproteinases.
  \textit{J Biol Chem} 274(31):21491--4.
  \textbf{nagase\_woessner1999.pdf} -- MMP overview. Contains no kinetic constants
  for TIMP-1/MMP-9; cited for biological context only.

\item Olson MW, Gervasi DC, Mobashery S, Fridman R (1997) Kinetic analysis of the
  binding of human matrix metalloproteinase-2 and -9 to tissue inhibitor of
  metalloproteinase (TIMP)-1 and TIMP-2.
  \textit{J Biol Chem} 272(47):29975--83.
  \textbf{olson1997.pdf} -- $\kon$, $\koff$, $K_i$ for MMP-9 (82~kDa)/TIMP-1
  confirmed from Table~IV (enzyme inhibition assay).

\item Beatty K, Bieth J, Travis J (1980) Kinetics of association of serine proteinases
  with native and oxidized $\alpha$-1-proteinase inhibitor and $\alpha$-1-antichymotrypsin.
  \textit{J Biol Chem} 255(9):3931--34.
  \textbf{beatty1980.pdf} -- $\kon$(HNE/$\alpha_1$-PI) $= 6.5\times10^7$~M$^{-1}$s$^{-1}$
  confirmed from Table~1.

\item Ashcroft GS, Lei K, Jin W et al.\ (2000) Secretory leukocyte protease inhibitor
  mediates non-redundant functions necessary for normal wound healing.
  \textit{Nat Med} 6(10):1147--53.
  \textbf{ashcroft2000.pdf} -- confirms SLPI role in wound healing (knockout mouse model).
  Does not provide SLPI concentration in human wound fluid; SLPI concentration remains
  unconfirmed in this model.

\item Rao CN, Ladin DA, Liu YY et al.\ (1995) $\alpha_1$-Antitrypsin is degraded and
  non-functional in chronic wounds but intact and functional in acute wounds: the inhibitor
  protects fibronectin from degradation by chronic wound fluid enzymes.
  \textit{J Invest Dermatol} 105(4):572--578.
  \textbf{rao1995.pdf} -- A1AT intact in all healing/acute samples; degraded and
  non-functional in 9/10 chronic wound samples. Confirmed from text (p.572--575).
  No quantitative A1AT concentration in ng/mL is provided; healing-state concentration
  remains an estimate (serum $\div$10).

\item Nwomeh BC, Liang HX, Cohen IK, Yager DR (1999) MMP-8 is the predominant collagenase
  in healing wounds and nonhealing ulcers.
  \textit{J Surg Res} 81(2):189--195.
  \textbf{nwomeh1999.pdf} -- TIMP-1 levels ``very low or not detected'' in chronic ulcer
  fluid (text p.193), confirming TIMP-1 suppression in chronic wound environments.
  Quantitative TIMP-1 values are in ng per 100~$\mu$g protein (not ng/mL); cannot be
  converted to absolute concentration without protein content data.
\end{enumerate}

\end{document}
